\documentclass[11pt,reqno,a4paper]{amsart}
\usepackage[a4paper,left=2.6cm,right=2.6cm,top=2.7cm,bottom=2.7cm]{geometry}
\usepackage{amsmath,amssymb,amsfonts}
\usepackage{array,booktabs,float}
\numberwithin{equation}{section}
\usepackage{hyperref}
\hypersetup{colorlinks=true,linkcolor=blue,citecolor=blue,urlcolor=blue}
\newtheorem{theorem}{Theorem}[section]
\newtheorem{lemma}[theorem]{Lemma}
\newtheorem{proposition}[theorem]{Proposition}
\newtheorem{corollary}[theorem]{Corollary}
\theoremstyle{definition}
\newtheorem{example}[theorem]{Example}
\newtheorem{remark}[theorem]{Remark}
\theoremstyle{plain}
\newtheorem*{theoremA}{Theorem A}
\newtheorem*{conjecture}{Conjecture}
\newcommand{\cf}{\mathrm{cf}}
\newcommand{\GL}{\mathrm{GL}}
\newcommand{\End}{\operatorname{End}}
\newcommand{\Imop}{\operatorname{Im}}
\newcommand{\id}{\mathrm{id}}
\newcommand{\ord}{\operatorname{ord}}
\newcommand{\ch}{\operatorname{char}}
\newcommand{\tr}{\operatorname{tr}}
\newcommand{\LT}{\mathrm{LT}_{\infty}}
\newcommand{\UT}{\mathrm{UT}_{\infty}}
\newcommand{\Mat}{M_{\cf}}
\newcommand{\bbN}{\mathbb N}
\newcommand{\bbZ}{\mathbb Z}
\newcommand{\bbC}{\mathbb C}
\newcommand{\bbF}{\mathbb F}
\newcommand{\spann}[1]{\langle #1\rangle}
\newcommand{\Acal}{\mathcal A}
\newcommand{\Jcal}{\mathcal J}
\newcommand{\Tcal}{\mathcal T}
\newcommand{\Scal}{\mathcal S}
\title[LOWER TRIANGULAR MATRICES AND PRODUCTS OF TWO INVOLUTIONS]{INFINITE LOWER TRIANGULAR MATRICES THAT ARE SIMILAR TO THEIR INVERSES BUT ARE NOT PRODUCTS OF TWO INVOLUTIONS}
\author[PHAM THE NHAN]{PHAM THE NHAN}
\address{Faculty of Mathematics and Computer Science, University of Science (HCMUS), Ho Chi Minh City, Vietnam}
\email{ptnhan@hcmus.edu.vn}
\subjclass[2020]{15A23, 20E45, 20H20}
\keywords{Infinite matrix; column-finite matrix; triangular matrix; involution; reversible element; Laurent polynomial}
% AUTHOR PRE-SUBMISSION CHECKS (not printed):
% 1. Verify and add exact AI product/model version numbers to the AI declaration;
%    check applicable terms of use and ensure compliance with T&F AI policy.
% 2. Independent research: no external research funding declared.
%    If the portal asks about funding, select no external funding.
% 3. Confirm that the sole-authorship, conflict-of-interest, and data statements
%    are correct, and independently verify all mathematical arguments.
% 4. Confirm originality and that this manuscript is not under review elsewhere.
% 5. Revisit journal-specific Instructions for Authors at the submission portal.
\begin{document}
\raggedbottom
\begin{abstract}
Let $F$ be a field, and let $\GL_{\cf}(F)$ be the group of invertible column-finite matrices over $F$ indexed by $\bbN\times\bbN$. Denote by $\LT(F)$ (resp.\ $\UT(F)$) the subgroup consisting of lower (resp.\ upper) triangular matrices whose diagonal entries are equal to $1$. Every matrix in $\UT(F)$ is a product of two involutions. We show that, if $F$ has at least four elements, then $\LT(F)$ contains a matrix which is conjugate to its inverse by an invertible lower triangular matrix, but is not a product of two involutions in $\GL_{\cf}(F)$. The matrices can be chosen to be $4$-periodic. This gives a negative answer to a conjecture of de Seguins Pazzis.
\end{abstract}
\maketitle

\section{Introduction}
Recall that an element $g$ of a group $G$ is called an involution if $g^2=1$. If $g=ab$, where $a$ and $b$ are involutions of $G$, then $bgb^{-1}=ba=g^{-1}$, that is, $g$ is conjugate to its inverse. Following \cite{ofarrell2015}, an element of $G$ is called \emph{reversible} if it is conjugate to its inverse, and it is called \emph{strongly reversible} if it is conjugate to its inverse by an involution; the latter condition means exactly that the element is a product of two involutions. In general, a reversible element is not necessarily strongly reversible. For example, every element of $\mathrm{SL}_2(\bbC)$ is conjugate to its inverse in $\mathrm{SL}_2(\bbC)$, while the only products of two involutions in $\mathrm{SL}_2(\bbC)$ are $\pm I_2$ (see \cite{pazzis2024}).

In the general linear group of a finite-dimensional vector space, these two notions coincide. In 1966, Wonenburger \cite{wonenburger1966} showed that if $V$ is a finite-dimensional vector space over a field of characteristic different from $2$, then an element of $\GL(V)$ is a product of two involutions if and only if it is similar to its inverse. The assumption on the characteristic was removed by Djokovi\'{c} \cite{djokovic1967} and by Hoffman and Paige \cite{hoffman1971}.

Let $F$ be a field and let $V$ be a vector space over $F$ with countably infinite dimension. Fixing a basis of $V$, we identify $\End(V)$ with the algebra $M_{\cf}(F)$ of column-finite matrices indexed by $\bbN\times\bbN$, and $\GL(V)$ with its unit group $\GL_{\cf}(F)$ (see Section 2.1). We denote by $\UT(F)$ (resp.\ $\LT(F)$) the subgroup consisting of upper (resp.\ lower) triangular matrices with all diagonal entries equal to $1$. S\l{}owik \cite{slowik2013} showed that every upper triangular matrix in $\GL_{\cf}(F)$ with diagonal entries $\pm1$ is a product of at most five involutions. Related results for matrices over rings and infinite upper triangular quadratic matrices were obtained in \cite{hou2017,bien2024}. De Seguins Pazzis \cite{pazzis2021} proved that every element of $\GL(V)$ is a product of four involutions and characterized the products of three involutions. He also showed that a locally finite automorphism of $V$ is a product of two involutions if and only if it is similar to its inverse \cite[Theorem 1.5]{pazzis2022}, and that every matrix in $\UT(F)$ is a product of two involutions \cite[Theorem 1.9]{pazzis2022}. He proposed the following conjecture.

\begin{conjecture}[{\cite[Conjecture 2]{pazzis2022}}]
Let $V$ be an infinite-dimensional vector space, and let $u\in\GL(V)$. Then, $u$ is the product of two involutions in $\GL(V)$ if and only if $u$ is similar to its inverse.
\end{conjecture}

For infinite matrices, lower and upper triangular matrices behave differently up to similarity. Indeed, every matrix in $\UT(F)$ differs from the identity by a locally nilpotent operator, and hence is locally finite. In contrast, $\LT(F)$ contains matrices that are not locally finite (Example 2.3), as well as matrices that are not similar to their inverses (Example 2.5). These distinctions are summarized in Table 1 and motivate testing the above conjecture in $\LT(F)$.

\begin{table}[ht]
\centering
\caption{Matrices in $\UT(F)$ and in $\LT(F)$.}
\label{tab:comparison}
\small
\begin{tabular}{@{}p{0.32\textwidth}p{0.27\textwidth}p{0.34\textwidth}@{}}
Property of $A$ & $A\in\UT(F)$ & $A\in\LT(F)$ \\[2pt]
$A-I$ is locally nilpotent & always & not always (Example 2.3) \\[2pt]
$A$ is similar to $A^{-1}$ & always & not always (Example 2.5) \\[2pt]
$A$ is a product of two involutions & always \cite[Theorem 1.9]{pazzis2022} & not always, even if $A$ is similar to $A^{-1}$ (Theorem A)
\end{tabular}
\end{table}

Our main result gives a negative answer to the conjecture.

\begin{theoremA}
Let $F$ be a field with at least four elements. Then, there exist a matrix $A=(a_{ij})\in\LT(F)$ and an invertible lower triangular matrix $P=(p_{ij})\in\GL_{\cf}(F)$ such that $PAP^{-1}=A^{-1}$, but $A$ is not a product of two involutions in $\GL_{\cf}(F)$. Moreover, $A$ and $P$ can be chosen to be $4$-periodic, that is, $a_{i+4,j+4}=a_{ij}$ and $p_{i+4,j+4}=p_{ij}$ for all $i,j$.
\end{theoremA}

The matrix $A$ in Theorem A comes from an explicit $F$-linear automorphism of the Laurent polynomial space $F[t,t^{-1}]$, while $P$ represents multiplication by $t$. The construction is given in Section 3, and Table 2 provides the corresponding matrices when $\ch F\ne2,3$. Theorem A applies over every field with at least four elements, including all algebraically closed fields, but leaves the cases $\bbF_2$ and $\bbF_3$ open.

Thus, the locally finite hypothesis in \cite[Theorem 1.5]{pazzis2022} cannot be omitted. A related distinction between reversible and strongly reversible elements was observed in the Riordan group \cite{slowik2026}, whose matrices are not column-finite in general.

This paper is organized as follows. In Section 2, we recall some facts on column-finite matrices, and we give examples illustrating the differences between $\LT(F)$ and $\UT(F)$. In Section 3, we construct the matrix $A$ in Theorem A, and we prove Theorem A.

Our notations in this paper are standard. Throughout this paper, $F$ is a field, $\bbN=\{1,2,\ldots\}$, and $V$ is a vector space over $F$ with countably infinite dimension. For a commutative ring $S$, we write $S^*$ for the unit group of $S$. We denote by $\id$ the identity map of $V$ (or of any other space under consideration) and by $I$ the identity matrix in $M_{\cf}(F)$. For $u\in\End(V)$, we denote by $\Imop u$ the image of $u$. For a set $X$ of vectors, we denote by $\spann{X}$ the subspace spanned by $X$.

\section{The groups \texorpdfstring{$\LT(F)$ and $\UT(F)$}{LT-infinity(F) and UT-infinity(F)}}
In this section, we compare the groups $\LT(F)$ and $\UT(F)$ up to similarity. We begin with some facts on column-finite matrices.

\subsection{Column-finite matrices}
Let $(e_n)_{n\ge1}$ be a basis of $V$. For $u\in\End(V)$, we write $u(e_j)=\sum_{i\ge1}a_{ij}e_i$, where, for each $j$, only finitely many $a_{ij}$ are nonzero. The matrix $(a_{ij})_{i,j\ge1}$ is called the matrix of $u$ in the basis $(e_n)$. The column-finite matrices over $F$ form an $F$-algebra $M_{\cf}(F)$, and the map sending $u$ to its matrix is an isomorphism of $F$-algebras $\End(V)\to M_{\cf}(F)$, which maps $\GL(V)$ onto the unit group $\GL_{\cf}(F)$ of $M_{\cf}(F)$. A change of basis replaces the matrix of $u$ by a similar matrix in $\GL_{\cf}(F)$, and two matrices in $M_{\cf}(F)$ are similar in $\GL_{\cf}(F)$ if and only if they are the matrices of the same endomorphism of $V$ in two bases. Therefore, the conjecture in Section 1 and our results can be stated equivalently for automorphisms of $V$ or for matrices in $\GL_{\cf}(F)$, and we will pass freely between these two points of view. In particular, if $W$ is a vector space with countably infinite dimension, then, after choosing a basis of $W$, we identify $\GL(W)$ with $\GL_{\cf}(F)$.

For $n\ge1$, put $V_n=\spann{e_m:m\ge n}$. The matrix of $u$ is lower triangular if and only if $u(V_n)\subseteq V_n$ for all $n$, and it is strictly lower triangular if and only if $u(V_n)\subseteq V_{n+1}$ for all $n$. Similarly, the matrix of $u$ is upper triangular if and only if $u(\spann{e_1,\ldots,e_n})\subseteq\spann{e_1,\ldots,e_n}$ for all $n$.

Following \cite{pazzis2022}, an endomorphism $u$ of $V$ is called locally nilpotent if, for every $x\in V$, there exists $n\ge1$ such that $u^n(x)=0$, and $u$ is called locally finite if, for every $x\in V$, the family $(u^k(x))_{k\ge0}$ is linearly dependent, that is, every $x\in V$ lies in a finite-dimensional $u$-invariant subspace. It is well known that $u$ is locally nilpotent if and only if the matrix of $u$ in some basis of $V$ is strictly upper triangular (see \cite[Section 1.2]{pazzis2022}). Consequently, an automorphism $A$ of $V$ is similar to a matrix in $\UT(F)$ if and only if $A-\id$ is locally nilpotent; in this case, $A$ is locally finite.

If $Y\in M_{\cf}(F)$ is an upper triangular matrix whose diagonal entries are equal to $1$, then $Y-I$ is locally nilpotent, so $Y$ is invertible in $M_{\cf}(F)$. This is not true for lower triangular matrices. For example, let $S\in M_{\cf}(F)$ be the matrix of the map $e_n\mapsto e_{n+1}$. Then, $I+S$ is lower triangular with diagonal entries $1$, but $e_1$ does not belong to the image of $I+S$. This is the reason why we require that the elements of $\LT(F)$ are invertible. The following lemma is elementary; for completeness, we provide a brief proof.

\begin{lemma}
Let $X\in\GL_{\cf}(F)$ be a lower triangular matrix. Then, $X^{-1}$ is lower triangular, and the diagonal entries of $X^{-1}$ are the inverses of those of $X$. In particular, $\LT(F)$ is a subgroup of $\GL_{\cf}(F)$.
\end{lemma}
\begin{proof}
Put $Y=X^{-1}$. For each $n\ge1$, let $X_n$ and $Y_n$ denote the leading principal submatrices of order $n$ of $X$ and $Y$, respectively. Since $X$ is lower triangular and $XY=I$, we have $X_nY_n=I_n$. Hence $Y_n=X_n^{-1}$, which is lower triangular with diagonal entries $x_{11}^{-1},\ldots,x_{nn}^{-1}$. As this holds for every $n$, the conclusion follows.
\end{proof}

\subsection{An adapted basis lemma}
We begin with the following lemma, which will be used frequently in this paper.

\begin{lemma}
Let $X_1=V\supseteq X_2\supseteq\cdots$ be subspaces of $V$ with $\bigcap_{n\ge1}X_n=0$. For $0\ne v\in V$, put $\ell(v)=\max\{n\ge1:v\in X_n\}$.
\begin{enumerate}
\item[(a)] Let $v_1,\ldots,v_r\in V$ be nonzero vectors such that $\ell(v_1),\ldots,\ell(v_r)$ are pairwise distinct, and let $c_1,\ldots,c_r\in F$ be not all zero. Then, $\sum_i c_iv_i\ne0$ and $\ell(\sum_i c_iv_i)=\min\{\ell(v_i):c_i\ne0\}$. In particular, $v_1,\ldots,v_r$ are linearly independent.
\item[(b)] Assume, in addition, that $\dim X_n/X_{n+1}=1$ for all $n\ge1$. Then, $V$ has a basis $(v_n)_{n\ge1}$ such that $X_n=\spann{v_m:m\ge n}$ for all $n\ge1$.
\end{enumerate}
\end{lemma}
\begin{proof}
(a) Let $n=\min\{\ell(v_i):c_i\ne0\}$, and let $j$ be the unique index such that $c_j\ne0$ and $\ell(v_j)=n$. Then, $\sum_i c_iv_i\equiv c_jv_j\pmod{X_{n+1}}$ and $c_jv_j\in X_n\setminus X_{n+1}$. Hence, $\sum_i c_iv_i\in X_n\setminus X_{n+1}$, that is, $\ell(\sum_i c_iv_i)=n$.

(b) Let $(w_k)_{k\ge1}$ be a basis of $V$, and put $E_k=\spann{w_1,\ldots,w_k}$ and $E_0=0$. Since $E_k$ is finite-dimensional and $\bigcap_{n\ge1}X_n=0$, there exists $N_k$ such that $E_k\cap X_{N_k}=0$. We choose $u_k\in w_k+E_{k-1}$ so that $\ell(u_k)$ is as large as possible. If $\ell(u_k)=\ell(u_i)=n$ for some $i<k$, then, since $\dim X_n/X_{n+1}=1$, there exists $\lambda\in F$ such that $u_k-\lambda u_i\in X_{n+1}$. This contradicts the choice of $u_k$, since $u_k-\lambda u_i\in w_k+E_{k-1}$. Hence, the levels $\ell(u_k)$ are pairwise distinct. Moreover, $\spann{u_1,\ldots,u_k}=E_k$ for all $k$, so $(u_k)_{k\ge1}$ is a basis of $V$.

Let $n\ge1$ and $v\in X_n\setminus X_{n+1}$. Writing $v$ as a linear combination of the $u_k$ and using (a), we see that $n=\ell(u_k)$ for some $k$. Let $v_n$ be the vector $u_k$ such that $\ell(u_k)=n$. Then, $(v_n)_{n\ge1}$ is a basis of $V$, and by (a), a vector $\sum_m c_mv_m$ belongs to $X_n$ if and only if $c_m=0$ for all $m<n$. Therefore, $X_n=\spann{v_m:m\ge n}$. The proof is complete.
\end{proof}

\subsection{Examples}
The following three examples show the differences between $\LT(F)$ and $\UT(F)$ in Table 1.

\begin{example}[A bilateral shift]
Let $f$ be the automorphism of $V$ defined by
\[
f(e_1)=e_2,\qquad f(e_{2n})=e_{2n+2},\qquad f(e_{2n+1})=e_{2n-1}\quad(n\ge1).
\]
Then, $f^k(e_1)=e_{2k}$ and $f^{-k}(e_1)=e_{2k+1}$ for all $k\ge1$, that is, $f$ is the bilateral shift along $\ldots,e_5,e_3,e_1,e_2,e_4,\ldots$. Put $g_1=f-\id$ and $g_2=\id-f^{-1}$. Then, $g_1$ and $g_2$ commute with $f$, and
\begin{equation}\label{eq:g12}
g_1=fg_2.
\end{equation}
We put $v_1=e_1$ and
\[
v_{2k}=g_1(v_{2k-1}),\qquad v_{2k+1}=g_2(v_{2k})\quad(k\ge1).
\]
For example, $v_2=e_2-e_1$, $v_3=e_3+e_2-2e_1$ and $v_4=e_4-e_3-3e_2+3e_1$.

First of all, we claim that $(v_n)_{n\ge1}$ is a basis of $V$. Since $g_2=f^{-1}g_1$ and $f$ commutes with $g_1$, induction on $k$ shows that
\[
v_{2k+1}=f^{-k}(f-\id)^{2k}(e_1),\qquad
v_{2k+2}=f^{-k}(f-\id)^{2k+1}(e_1)\quad(k\ge0).
\]
By the binomial formula, the term in $v_{2k+1}$ involving $f^{-k}(e_1)=e_{2k+1}$ has coefficient $1$, while all other terms belong to $\spann{e_1,\ldots,e_{2k}}$. Likewise, the term in $v_{2k+2}$ involving $f^{k+1}(e_1)=e_{2k+2}$ has coefficient $1$, while all other terms belong to $\spann{e_1,\ldots,e_{2k+1}}$. Therefore,
\[
v_n-e_n\in\spann{e_1,\ldots,e_{n-1}}\qquad\text{for all }n\ge1.
\]
It follows that $\spann{v_1,\ldots,v_n}=\spann{e_1,\ldots,e_n}$ for all $n$, so $(v_n)_{n\ge1}$ is a basis of $V$.

Next, we compute the matrix of $f$ in the basis $(v_n)$. By the definitions of $g_1,g_2$ and \eqref{eq:g12},
\[
\begin{aligned}
f(v_{2k+1})&=v_{2k+1}+v_{2k+2},\\
f(v_{2k+2})&=v_{2k+2}+v_{2k+3}+v_{2k+4}\qquad(k\ge0).
\end{aligned}
\]
Therefore, the matrix of $f$ in the basis $(v_n)$ is
\[
A=\begin{pmatrix}
1&0&0&0&0&0&\cdots\\
1&1&0&0&0&0&\cdots\\
0&1&1&0&0&0&\cdots\\
0&1&1&1&0&0&\cdots\\
0&0&0&1&1&0&\cdots\\
0&0&0&1&1&1&\cdots\\
\vdots&\vdots&\vdots&\vdots&\vdots&\vdots&\ddots
\end{pmatrix}\in\LT(F).
\]
On the other hand, $f$ is not locally finite, since the family $(f^k(e_1))_{k\ge0}=(e_1,e_2,e_4,\ldots)$ is linearly independent. Hence $f-\id$ is not locally nilpotent, and by Section 2.1, $A$ is not similar to any matrix in $\UT(F)$.

However, $f$ is a product of two involutions. Indeed, let $S$ and $T$ be the automorphisms of $V$ defined by $S(f^k(e_1))=f^{-k}(e_1)$ and $T(f^k(e_1))=f^{1-k}(e_1)$ for $k\in\bbZ$, that is, $S$ fixes $e_1$ and interchanges $e_{2k}$ and $e_{2k+1}$ for every $k\ge1$, and $T$ interchanges $e_{2k-1}$ and $e_{2k}$ for every $k\ge1$. Then, $S^2=T^2=\id$, and $TS(f^k(e_1))=T(f^{-k}(e_1))=f^{k+1}(e_1)$ for all $k\in\bbZ$, that is, $f=TS$.
\end{example}

\begin{example}
Let $B\in\UT(F)$ be the matrix such that $B(e_1)=e_1$ and $B(e_{n+1})=e_{n+1}+e_n$ for $n\ge1$. Then, $B-I$ maps $e_{n+1}$ to $e_n$, so $\Imop(B-I)^n=V$ for every $n$. On the other hand, for every $C\in\LT(F)$, we have
\[
\Imop(C-I)^n\subseteq\spann{e_j:j\ge n+1},
\]
and hence $\bigcap_{n\ge1}\Imop(C-I)^n=0$. Thus, $B$ is not similar to any matrix in $\LT(F)$.
\end{example}

\begin{example}[A matrix in $\LT(F)$ that is not similar to its inverse]
Let $f\in F[t]$ be an irreducible polynomial such that $f(0)\ne0$ and $f$ is not associate to its reciprocal polynomial $f^{\sharp}(t)=t^{\deg f}f(t^{-1})$. Such a polynomial exists over every field. For example, if $|F|>3$, then we can take $f=t-a$, where $a\in F\setminus\{0,1,-1\}$; if $F=\bbF_2$, then we can take $f=t^3+t+1$; and if $F=\bbF_3$, then we can take $f=t^2+t+2$. Consider the ring $M=F[t,t^{-1},f^{-1}]\subseteq F(t)$ and the automorphism $A$ of $M$ given by $A(h)=th$. Since $t-1$ is associate to its reciprocal polynomial $1-t$, the polynomial $f$ is not associate to $t-1$. Hence, every element of $M$ has nonnegative $(t-1)$-adic valuation, and so $\bigcap_{n\ge1}(t-1)^nM=0$.

Put $X_n=(t-1)^{n-1}M$ for $n\ge1$. Since evaluation at $t=1$ induces an isomorphism $M/(t-1)M\simeq F$, we have $\dim_F X_n/X_{n+1}=1$. By Lemma 2.2(b), there is a basis $(e_n)_{n\ge1}$ of $M$ such that $X_n=\spann{e_j:j\ge n}$. Since $(A-\id)(X_n)=X_{n+1}$, the matrix of $A$ in this basis belongs to $\LT(F)$.

Now, $f(A)$ is the multiplication by the unit $f$ of $M$, so it is bijective. On the other hand, $f(A^{-1})$ is the multiplication by $f(t^{-1})=t^{-\deg f}f^{\sharp}(t)$, and $f^{\sharp}$ is an irreducible polynomial which is not associate to $t$ or $f$. Hence, $f^{\sharp}$ is not a unit of $M$, and so $f(A^{-1})$ is not surjective. If $XAX^{-1}=A^{-1}$ for some $X\in\GL(M)$, then $Xf(A)X^{-1}=f(A^{-1})$, which is impossible. Therefore, $A$ is not similar to $A^{-1}$. In particular, $A$ is not a product of two involutions.
\end{example}

\section{A reversible matrix in \texorpdfstring{$\LT(F)$}{LT-infinity(F)} that is not a product of two involutions}
Throughout this section, let $F$ be a field with at least four elements. Fix $u\in F\setminus\{0,1,-1\}$, and put $m=u-1$.

\subsection{The construction}
Let $V=F[t,t^{-1}]$ be the ring of Laurent polynomials in $t$ over $F$. Then, $V$ is a vector space with basis $(t^n)_{n\in\bbZ}$, so $V$ has countably infinite dimension. For $g\in V$, we use the same letter $g$ for the map $V\to V$, $f\mapsto gf$; if $g$ is a unit of $V$, for example, $g=t$, then $g\in\GL(V)$. We begin with the following lemma.

\begin{lemma}
Let $g\in V$ be a Laurent polynomial with at least two nonzero coefficients. Then, $\bigcap_{n\ge1}g^nV=0$.
\end{lemma}
\begin{proof}
For $0\ne f\in V$, denote by $\delta(f)$ the difference between the largest and the smallest exponents of $t$ appearing in $f$. Since $F$ is a field, $\delta(ff')=\delta(f)+\delta(f')$ for all nonzero $f,f'\in V$. By assumption, $\delta(g)\ge1$, so $\delta(f)\ge n$ for every nonzero $f\in g^nV$. Hence, a nonzero element $f$ of $V$ does not belong to $g^nV$ whenever $n>\delta(f)$. The proof is complete.
\end{proof}

Put $\sigma=t^2$ and $\Lambda=F[\sigma,\sigma^{-1}]\subseteq V$. For $0\ne\gamma\in\Lambda$, we denote by $\deg\gamma$ and $\ord\gamma$ the largest and the smallest exponents of $\sigma$ appearing in $\gamma$, respectively, and we put $\ord0=\infty$. Since $F$ is a field, $\deg(\gamma\gamma')=\deg\gamma+\deg\gamma'$ and $\ord(\gamma\gamma')=\ord\gamma+\ord\gamma'$ for all nonzero $\gamma,\gamma'\in\Lambda$. In particular, $\Lambda$ is an integral domain, and
\[
\Lambda^*=\{\lambda\sigma^n:\lambda\in F^*,\ n\in\bbZ\}.
\]
Indeed, if $\gamma\gamma'=1$, then $\deg\gamma+\deg\gamma'=0=\ord\gamma+\ord\gamma'$, and since $\ord\gamma\le\deg\gamma$ and $\ord\gamma'\le\deg\gamma'$, we get $\ord\gamma=\deg\gamma$.

Every $f\in V$ can be written uniquely as $f=f_0+tf_1$, with $f_0,f_1\in\Lambda$. Thus, $V$ is a free $\Lambda$-module of rank two with basis $1,t$. Consequently,
\[
\End_{\Lambda}(V)\simeq M_2(\Lambda),
\]
and a $\Lambda$-linear map is bijective if and only if the determinant of its matrix belongs to $\Lambda^*$. Moreover, an $F$-linear map $T:V\to V$ is $\Lambda$-linear if and only if $T\sigma=\sigma T$.

Under this identification, multiplication by $\gamma\in\Lambda$ has matrix $\gamma I$, while multiplication by $t$ has matrix
\[
\Jcal=\begin{pmatrix}0&\sigma\\1&0\end{pmatrix}.
\]

Now put
\[
p=1+(2u-1)\sigma+u^2\sigma^2,\qquad r=m+u^2\sigma,\qquad s=1-u^2\sigma,
\]
and let $A:V\to V$ be the $\Lambda$-linear map with matrix
\[
\Acal=\begin{pmatrix}p&-\sigma r\\r&s\end{pmatrix}.
\]
Since $A(1)=p+tr$ and $A(t)=-\sigma r+ts=t(s-tr)$, we have
\begin{equation}\label{eq:Aaction}
A(f_0+tf_1)=af_0+ta'f_1\qquad(f_0,f_1\in\Lambda),
\end{equation}
where
\[
\begin{aligned}
a&=p+tr=1+mt+(2u-1)t^2+u^2t^3+u^2t^4,\\
a'&=s-tr=1-mt-u^2t^2-u^2t^3.
\end{aligned}
\]
We also put
\[
q=p+s=2-m^2\sigma+u^2\sigma^2\qquad\text{and}\qquad w=u^2\sigma-m^2,
\]
so that $q-2=\sigma w$.

\begin{lemma}
The following statements hold.
\begin{enumerate}
\item[(a)] $\det\Acal=1$. Hence, $A\in\GL(V)$, and the matrix of $A^{-1}$ is $\begin{pmatrix}s&\sigma r\\-r&p\end{pmatrix}$.
\item[(b)] $A+A^{-1}$ is the multiplication by $q$, and $(A-\id)^2=\sigma w A$.
\item[(c)] $tAt^{-1}=A^{-1}$.
\end{enumerate}
\end{lemma}
\begin{proof}
(a) Since $m=u-1$, a direct calculation gives
\[
\det\Acal=ps+\sigma r^2
=\bigl(1-m^2\sigma-2mu^2\sigma^2-u^4\sigma^3\bigr)
 +\bigl(m^2\sigma+2mu^2\sigma^2+u^4\sigma^3\bigr)=1.
\]
Hence, $A^{-1}$ is the adjugate matrix of $\Acal$, as stated.

(b) By (a), the matrix of $A+A^{-1}$ is $(p+s)I=qI$. Since $A$ commutes with the multiplication by $q\in\Lambda$, we get $A^2+\id=qA$, and so $(A-\id)^2=A^2-2A+\id=(q-2)A=\sigma wA$.

(c) By direct calculation,
\[
\Jcal\Acal=\begin{pmatrix}\sigma r&\sigma s\\p&-\sigma r\end{pmatrix}=\Acal^{-1}\Jcal,
\]
that is, $tA=A^{-1}t$. The proof is complete.
\end{proof}

\subsection{The matrix of \texorpdfstring{$A$}{A} in a suitable basis}
Put $\xi=-m/u$. Then, $\xi\ne0$, $ut+m=u(t-\xi)$ and $u^2\xi^2=m^2$, so that
\[
w=u^2(t-\xi)(t+\xi).
\]
For $n\ge0$, put $W_n=\Imop(A-\id)^n$; in particular, $W_0=V$.

\begin{lemma}
The following statements hold.
\begin{enumerate}
\item[(a)] $W_{2k}=w^kV$ and $W_{2k+1}=w^k(t-\xi)V$ for all $k\ge0$.
\item[(b)] $\dim W_n/W_{n+1}=1$ and $t(W_n)=W_n$ for all $n\ge0$, and $\bigcap_{n\ge0}W_n=0$.
\end{enumerate}
\end{lemma}
\begin{proof}
(a) First of all, we claim that $W_1=(t-\xi)V$. By direct calculation,
$a-1=ut(t-\xi)(1+t+ut^2)$ and $a'-1=-ut(t-\xi)(1+ut)$. Hence, by \eqref{eq:Aaction}, for $f_0,f_1\in\Lambda$, we have
\[
(A-\id)(f_0+tf_1)=(a-1)f_0+t(a'-1)f_1=ut(t-\xi)C(f_0+tf_1),
\]
where $C$ is the $\Lambda$-linear map given by $C(f_0+tf_1)=(1+t+ut^2)f_0-t(1+ut)f_1$. Since $C(1)=(1+u\sigma)+t$ and $C(t)=-u\sigma-t$, the matrix of $C$ is
\[
\begin{pmatrix}1+u\sigma&-u\sigma\\1&-1\end{pmatrix},
\]
whose determinant is $-1$. Hence, $C$ is bijective, and $W_1=ut(t-\xi)C(V)=(t-\xi)V$, since $ut$ is a unit of $V$. The claim is shown.

Now, by Lemma 3.2(b), $(A-\id)^2=\sigma wA$. Since $\sigma$ and $A$ are bijective and $A$ commutes with the multiplication by $w\in\Lambda$, we get $W_{2k}=(\sigma wA)^k(V)=w^k\sigma^k A^k(V)=w^kV$. Hence, $W_{2k+1}=(A-\id)(w^kV)=w^k(A-\id)(V)=w^k(t-\xi)V$.

(b) For every $\zeta\in F^*$, evaluation at $\zeta$ yields
\[
V=F\oplus(t-\zeta)V.
\]
Applying this decomposition with $\zeta=\xi$, we obtain
\[
W_{2k}=w^kV=Fw^k\oplus w^k(t-\xi)V=Fw^k\oplus W_{2k+1}.
\]
Similarly, applying it with $\zeta=-\xi$ and using $w=u^2(t-\xi)(t+\xi)$, we get
\[
\begin{aligned}
W_{2k+1}&=w^k(t-\xi)V\\
&=Fw^k(t-\xi)\oplus w^k(t-\xi)(t+\xi)V\\
&=Fw^k(t-\xi)\oplus W_{2k+2}.
\end{aligned}
\]
Here we use the fact that $V$ is an integral domain and $u^2\in F^*$. Consequently, $\dim W_n/W_{n+1}=1$ for all $n\ge0$. Moreover, $t(W_n)=W_n$, since each $W_n$ is an ideal of $V$ and $t$ is a unit of $V$. Finally,
\[
\bigcap_{n\ge0}W_n\subseteq\bigcap_{k\ge0}w^kV=0
\]
by Lemma 3.1, since $w=u^2t^2-m^2$ has two nonzero terms. The proof is complete.
\end{proof}

\begin{lemma}
Put $h=t-\xi^2t^{-1}$ and
\[
e_{2k+1}=h^k,\qquad e_{2k+2}=h^k(t-\xi)\quad(k\ge0).
\]
Then, $(e_n)_{n\ge1}$ is a basis of $V$, and $W_{n-1}=\spann{e_j:j\ge n}$ for all $n\ge1$.
\end{lemma}
\begin{proof}
First of all, we show that $V$ is spanned by the vectors $e_n$. Put $M=F[h]+F[h]t$. Since $t^2=ht+\xi^2$ and $t^{-1}=\xi^{-2}(t-h)$, we have $tM\subseteq M$ and $t^{-1}M\subseteq M$. Since $1\in M$, it follows that $t^n\in M$ for all $n\in\bbZ$, so $M=V$. Since $F[h]t\subseteq F[h](t-\xi)+F[h]$, we get $V=F[h]+F[h](t-\xi)$, that is, $V$ is spanned by the vectors $e_n$.

Next, we show that $e_n\in W_{n-1}\setminus W_n$ for all $n\ge1$. Since $w=u^2th$ and $u^2t$ is a unit of $V$, Lemma 3.3(a) gives
\[
W_{2k}=h^kV,\qquad W_{2k+1}=h^k(t-\xi)V\quad(k\ge0).
\]
As $h=t^{-1}(t-\xi)(t+\xi)$ and neither $t-\xi$ nor $t+\xi$ is a unit of $V$, we have
\[
e_{2k+1}\in W_{2k}\setminus W_{2k+1},\qquad
e_{2k+2}\in W_{2k+1}\setminus W_{2k+2}.
\]
Now we apply Lemma 2.2(a) to the subspaces $X_n=W_{n-1}$ $(n\ge1)$, whose intersection is $0$ by Lemma 3.3(b). Since $e_n\in X_n\setminus X_{n+1}$ for all $n$, the vectors $e_n$ are linearly independent, so $(e_n)$ is a basis of $V$. Moreover, a vector $\sum_j c_je_j$ with $c_j\in F$ belongs to $X_n=W_{n-1}$ if and only if $c_j=0$ for all $j<n$, that is, $W_{n-1}=\spann{e_j:j\ge n}$. The proof is complete.
\end{proof}

\begin{proposition}
In the basis $(e_n)$, the matrix of $A$ belongs to $\LT(F)$, and the matrix of $t$ is lower triangular; more precisely,
\[
te_{2k+1}=\xi e_{2k+1}+e_{2k+2},\qquad
te_{2k+2}=-\xi e_{2k+2}+\xi e_{2k+3}+e_{2k+4}\quad(k\ge0).
\]
Moreover, the matrices of $A$ and $t$ in this basis are $4$-periodic, that is, if $Y\in\{A,t\}$ and $Y(e_j)=\sum_i y_{ij}e_i$ for all $j$, then $y_{i+4,j+4}=y_{ij}$ for all $i,j\ge1$.
\end{proposition}
\begin{proof}
Let $n\ge1$. By Lemma 3.4, $e_n\in W_{n-1}$, so $(A-\id)(e_n)\in W_n=\spann{e_j:j\ge n+1}$. Hence, the matrix of $A-\id$ is strictly lower triangular, that is, the matrix of $A$ belongs to $\LT(F)$. Next, $te_{2k+1}=h^kt=h^k(t-\xi)+\xi h^k=e_{2k+2}+\xi e_{2k+1}$. Since $t^2=ht+\xi^2$, we have $t^2-\xi t=h(t-\xi)+\xi h-\xi(t-\xi)$, and so
\[
te_{2k+2}=h^k(t^2-\xi t)=e_{2k+4}+\xi e_{2k+3}-\xi e_{2k+2}.
\]
Finally, $h^2=w^2u^{-4}\sigma^{-1}\in\Lambda$, and $h^2e_j=e_{j+4}$ for all $j\ge1$. Since $A$ and $t$ commute with the multiplication by $h^2$, we get $Y(e_{j+4})=h^2Y(e_j)=\sum_i y_{ij}e_{i+4}$ for $Y\in\{A,t\}$, that is, $y_{i+4,j+4}=y_{ij}$. The proof is complete.
\end{proof}

\subsection{The automorphism \texorpdfstring{$A$}{A} is not a product of two involutions}
We begin with the following lemma, which is the key point of this subsection.

\begin{lemma}
Let $S:V\to V$ be an $F$-linear map such that $SA=A^{-1}S$. Then, $S$ is $\Lambda$-linear.
\end{lemma}
\begin{proof}
Multiplying both sides of $SA=A^{-1}S$ by $A$ on the left and by $A^{-1}$ on the right, we get $AS=SA^{-1}$. Hence, $S$ commutes with $A+A^{-1}$, which is the multiplication by $q$ by Lemma 3.2(b). Since $2-q=-\sigma w=u^2\sigma(\xi^2-\sigma)$, the map $S$ commutes with the multiplication by $\sigma(\xi^2-\sigma)$. Put $D=S\sigma-\sigma S$. It follows that
\[
D\sigma=(\xi^2-\sigma)D.
\]
Thus, $D=(\xi^2-\sigma)D\sigma^{-1}$, and by induction, $D=(\xi^2-\sigma)^nD\sigma^{-n}$ for all $n\ge1$. Hence,
\[
D(V)\subseteq\bigcap_{n\ge1}(\xi^2-\sigma)^nV=0
\]
by Lemma 3.1, since $\xi^2-\sigma=\xi^2-t^2$ and $\xi\ne0$. Therefore, $S\sigma=\sigma S$, that is, $S$ is $\Lambda$-linear. The proof is complete.
\end{proof}

\begin{lemma}
Let $\Tcal\in M_2(\Lambda)$ be such that $\Tcal\Acal=\Acal\Tcal$. Then, $\Tcal=\alpha I+\beta\Acal$ for some $\alpha,\beta\in\Lambda$.
\end{lemma}
\begin{proof}
Write $\Tcal=\begin{pmatrix}x&y\\z&v\end{pmatrix}$. Comparing the $(1,1)$ entries and the $(2,1)$ entries of $\Tcal\Acal$ and $\Acal\Tcal$, we get
\[
xp+yr=px-\sigma rz\qquad\text{and}\qquad zp+vr=rx+sz,
\]
that is, $yr=-\sigma rz$ and $r(v-x)=(s-p)z$. Since $u\ne0$, we have $r\ne0$, and since $\Lambda$ is an integral domain, $y=-\sigma z$. By direct calculation, $s-p=-\sigma(r+u(u+1))$, so
\[
r(v-x+\sigma z)=-\sigma u(u+1)z.
\]
Since $u\ne0$ and $u+1\ne0$, we have $\sigma u(u+1)\in\Lambda^*$, and so $z=\beta r$ for some $\beta\in\Lambda$. Then, $r(v-x)=(s-p)\beta r$, so $v-x=(s-p)\beta$. Therefore,
\[
\Tcal=\begin{pmatrix}x&-\sigma\beta r\\\beta r&x+(s-p)\beta\end{pmatrix}
=(x-\beta p)I+\beta\Acal.
\]
The proof is complete.
\end{proof}

In the next two lemmas, we use the following elementary fact: if finitely many nonzero elements of $\Lambda$ are given and exactly one of them has the smallest $\ord$, then the $\ord$ of their sum is this smallest $\ord$; the same holds for $\deg$ and the largest $\deg$.

\begin{lemma}
Let $\alpha,\beta\in\Lambda$ with $\beta\ne0$. Then,
\[
\ord(\alpha^2+q\alpha\beta+\beta^2)=\min\{2\ord(\alpha+\beta),\ 2\ord\beta+1\}.
\]
\end{lemma}
\begin{proof}
Put $\gamma=\alpha+\beta$. Since $q-2=\sigma w$, we have
\[
\alpha^2+q\alpha\beta+\beta^2=(\gamma-\beta)^2+q(\gamma-\beta)\beta+\beta^2
=\gamma^2+\sigma w\gamma\beta-\sigma w\beta^2.
\]
Let $j=\ord\beta$ and $k=\ord\gamma$, with $k=\infty$ if $\gamma=0$. Since the constant term of $w$ is $-m^2\ne0$, we have $\ord w=0$. The three terms on the right have orders $2k$, $k+j+1$ and $2j+1$, respectively, whenever they are nonzero. Since $k+j+1>\min\{2k,2j+1\}$ and $2k\ne2j+1$, the smallest of these orders occurs exactly once. Therefore, the order of the sum is $\min\{2k,2j+1\}$. The proof is complete.
\end{proof}

\begin{lemma}
Let $\alpha,\beta\in\Lambda$, and assume that $\alpha^2+q\alpha\beta+\beta^2=\lambda\sigma^n$ for some $\lambda\in F^*$ and $n\in\bbZ$. Then, $n$ is even, and $\alpha I+\beta\Acal=\lambda'\sigma^{n/2}\Acal^i$ for some $\lambda'\in F^*$ and $i\in\bbZ$.
\end{lemma}
\begin{proof}
If we replace $\alpha$ and $\beta$ by $\sigma^j\alpha$ and $\sigma^j\beta$ with $j\in\bbZ$, then $n$ is replaced by $n+2j$, and $\alpha I+\beta\Acal$ is replaced by $\sigma^j(\alpha I+\beta\Acal)$. Hence, we may assume that $n\in\{0,1\}$. Put $\Tcal=\alpha I+\beta\Acal$. Recall that $\det(\alpha I+\beta M)=\alpha^2+\alpha\beta\tr M+\beta^2\det M$ for every $M\in M_2(\Lambda)$. Hence, $\det\Tcal=\alpha^2+q\alpha\beta+\beta^2=\lambda\sigma^n$.

By Lemma 3.2(b), we have $\Acal^2=q\Acal-I$. Hence, for every $i\in\bbZ$, we can write $\Acal^i\Tcal=\alpha_iI+\beta_i\Acal$ with $\alpha_i,\beta_i\in\Lambda$. These elements are uniquely determined, since the $(2,1)$ entry of $\alpha_iI+\beta_i\Acal$ is $\beta_i r$ and $r\ne0$. Since
\[
\Acal^{i+1}\Tcal=\Acal(\alpha_iI+\beta_i\Acal)=-\beta_iI+(\alpha_i+q\beta_i)\Acal,
\]
we get $\alpha_{i+1}=-\beta_i$ and $\beta_{i+1}=\alpha_i+q\beta_i$ for all $i\in\bbZ$, that is,
\begin{equation}\label{eq:recurrence}
\alpha_i=-\beta_{i-1}\qquad\text{and}\qquad\beta_{i-1}+\beta_{i+1}=q\beta_i\quad(i\in\bbZ).
\end{equation}
Moreover, $\det(\Acal^i\Tcal)=\det\Tcal=\lambda\sigma^n$, and by \eqref{eq:recurrence},
\[
\det(\alpha_iI+\beta_i\Acal)=\beta_{i-1}^2-q\beta_{i-1}\beta_i+\beta_i^2=\beta_i^2-\beta_{i-1}\beta_{i+1}.
\]
Hence,
\begin{equation}\label{eq:detrecurrence}
\beta_i^2-\beta_{i-1}\beta_{i+1}=\lambda\sigma^n\qquad(i\in\bbZ).
\end{equation}

\emph{Case 1: $\beta_i=0$ for some $i$.} Then, $\Acal^i\Tcal=\alpha_iI$ and $\alpha_i^2=\det(\alpha_iI)=\lambda\sigma^n$. Hence, $2\ord\alpha_i=n=2\deg\alpha_i$, so $n=0$ and $\alpha_i\in F^*$. Thus, $\Tcal=\alpha_i\Acal^{-i}$, as desired.

\emph{Case 2: $\beta_i\ne0$ for all $i$.} We show that this case does not occur. Let $i\in\bbZ$. Since $\alpha_i^2+q\alpha_i\beta_i+\beta_i^2=\det(\Acal^i\Tcal)=\lambda\sigma^n$, Lemma 3.8 applied to the pair $(\alpha_i,\beta_i)$ gives $n=\ord(\lambda\sigma^n)\le2\ord\beta_i+1$. Since $n\ge0$, we get $\ord\beta_i\ge0$. Thus, $\deg\beta_i\ge\ord\beta_i\ge0$ for all $i$, and we can choose $i$ such that $d=\deg\beta_i$ is the smallest. Since $u\ne0$, we have $\deg q=2$, so by \eqref{eq:recurrence}, $\beta_{i-1}+\beta_{i+1}=q\beta_i$ has degree $d+2$. Hence, one of $\beta_{i-1}$ and $\beta_{i+1}$ has degree at least $d+2$, and the other has degree at least $d$ by the choice of $i$. Thus, $\deg(\beta_{i-1}\beta_{i+1})\ge2d+2>\deg\beta_i^2$, and by \eqref{eq:detrecurrence},
\[
n=\deg(\lambda\sigma^n)=\deg(\beta_i^2-\beta_{i-1}\beta_{i+1})
=\deg(\beta_{i-1}\beta_{i+1})\ge2d+2\ge2,
\]
which contradicts $n\le1$. The proof is complete.
\end{proof}

Now we are ready to show the main result of this section.

\begin{theorem}
Let $S\in\GL(V)$ be such that $SAS^{-1}=A^{-1}$. Then, $S=\lambda\sigma^jA^it$ for some $\lambda\in F^*$ and $i,j\in\bbZ$, and $S^2$ is the multiplication by $\lambda^2t^{4j+2}$. In particular, $S^2\ne\id$, and $A$ is not a product of two involutions in $\GL(V)$.
\end{theorem}
\begin{proof}
By Lemma 3.6, $S$ is $\Lambda$-linear. Let $\Scal$ be the matrix of $S$. Then, $\Scal$ is invertible, since $S$ is bijective, and $\Scal\Acal\Scal^{-1}=\Acal^{-1}$. Put $\Tcal=\Scal\Jcal^{-1}$. By Lemma 3.2(c), both $\Scal$ and $\Jcal$ conjugate $\Acal$ to $\Acal^{-1}$, so $\Tcal$ commutes with $\Acal$. By Lemma 3.7, $\Tcal=\alpha I+\beta\Acal$ for some $\alpha,\beta\in\Lambda$. Since $\Tcal$ is invertible, $\det\Tcal=\alpha^2+q\alpha\beta+\beta^2$ belongs to $\Lambda^*$, so it is of the form $\lambda_0\sigma^n$ with $\lambda_0\in F^*$ and $n\in\bbZ$. By Lemma 3.9, $\Tcal=\lambda\sigma^j\Acal^i$ for some $\lambda\in F^*$ and $i,j\in\bbZ$. Hence, $\Scal=\lambda\sigma^j\Acal^i\Jcal$, that is, $S=\lambda\sigma^jA^it$. Since $\sigma$ commutes with $A$ and $t$, and $tA^it^{-1}=A^{-i}$ by Lemma 3.2(c), we get
\[
S^2=\lambda^2\sigma^{2j}A^i(tA^it^{-1})t^2
=\lambda^2\sigma^{2j}A^iA^{-i}\sigma=\lambda^2\sigma^{2j+1},
\]
which is the multiplication by $\lambda^2t^{4j+2}$. Since $4j+2\ne0$, we have $S^2\ne\id$.

Finally, suppose that $A=S_1S_2$ for some involutions $S_1,S_2\in\GL(V)$. Then, $S_1AS_1^{-1}=S_1S_1S_2S_1=S_2S_1=A^{-1}$, so $S_1^2\ne\id$ by the above, a contradiction. Therefore, $A$ is not a product of two involutions in $\GL(V)$. The proof is complete.
\end{proof}

\subsection{Proof of Theorem A}
\begin{proof}[Proof of Theorem A]
Let $F$ be a field with at least four elements. Then, there exists $u\in F$ such that $u\notin\{0,1,-1\}$, and we use the above notation. As in Section 2.1, we identify $\GL(V)$ with $\GL_{\cf}(F)$ via the basis $(e_n)$ of Lemma 3.4, and we use the same letter $A$ for the matrix of $A$ in this basis. Let $P$ be the matrix of the multiplication by $t$ in this basis. By Proposition 3.5, $A\in\LT(F)$, $P$ is an invertible lower triangular matrix, and both $A$ and $P$ are $4$-periodic. By Lemma 3.2(c), $PAP^{-1}=A^{-1}$. By Theorem 3.10, $A$ is not a product of two involutions in $\GL_{\cf}(F)$. The proof is complete.
\end{proof}

We end this paper with an explicit example.

\begin{example}
Assume that $\ch F\ne2,3$, and take $u=2$. Then, $m=1$, $\xi=-\frac12$, $h=t-\frac14t^{-1}$, and by \eqref{eq:Aaction},
\[
A(f_0+tf_1)=(1+t+3t^2+4t^3+4t^4)f_0+t(1-t-4t^2-4t^3)f_1\quad(f_0,f_1\in\Lambda).
\]
By Proposition 3.5, the matrices of $A$ and $2t$ in the basis $(e_n)$ are $4$-periodic, so they are determined by their first four columns. A direct calculation gives Table 2. For example, the first row of Table 2 with $j=1$ means that
\[
1+t+3t^2+4t^3+4t^4=e_1+2e_2-\tfrac32e_3+5e_4-e_5+4e_6-2e_7+4e_8.
\]
\end{example}

\begin{table}[H]
\centering
\caption{The $4$-periodic matrices of $A$ and $2t$ in Example 3.11: the images of $e_j$ according to the residue of $j$ modulo $4$.}
\label{tab:explicit}
\small
\renewcommand{\arraystretch}{1.3}
\begin{tabular}{@{}c@{\quad}p{0.83\textwidth}@{}}
$j\bmod4$ & $A(e_j)$\\\hline
$1$ & $e_j+2e_{j+1}-\tfrac32e_{j+2}+5e_{j+3}-e_{j+4}+4e_{j+5}-2e_{j+6}+4e_{j+7}$\\
$2$ & $e_j-\tfrac14e_{j+1}-\tfrac12e_{j+2}+\tfrac12e_{j+3}-2e_{j+4}+e_{j+5}-2e_{j+6}$\\
$3$ & $e_j-2e_{j+1}+e_{j+2}-4e_{j+3}+2e_{j+4}-4e_{j+5}$\\
$0$ & $e_j+\tfrac14e_{j+1}+e_{j+2}-e_{j+3}+4e_{j+4}-e_{j+5}+4e_{j+6}-2e_{j+7}+4e_{j+8}$\\[2pt]
$j\bmod2$ & $2t\,e_j$\\\hline
$1$ & $-e_j+2e_{j+1}$\\
$0$ & $e_j-e_{j+1}+2e_{j+2}$
\end{tabular}
\end{table}

\section*{Declaration of generative AI use}
During the preparation of this work, the authors used AI-assisted tools to support literature searches and improve the readability and language of the manuscript. After using these tools, the authors independently reviewed and edited all incorporated content and take full responsibility for the final published work.

\section*{Disclosure statement}
The author declares no competing interests.

\section*{Data availability statement}
No datasets were generated or analyzed during this study.

\begin{thebibliography}{11}
\bibitem{bien2024} M. H. Bien, V. M. Tam, D. C. M. Tri, and L. Q. Truong, Products of infinite upper triangular quadratic matrices, \emph{Linear Algebra Appl.} \textbf{699} (2024) 59--71.
\bibitem{djokovic1967} D. \v{Z}. Djokovi\'{c}, Product of two involutions, \emph{Arch. Math. (Basel)} \textbf{18} (1967) 582--584.
\bibitem{hoffman1971} F. Hoffman and E. C. Paige, Products of two involutions in the general linear group, \emph{Indiana Univ. Math. J.} \textbf{20} (1971) 1017--1020.
\bibitem{hou2017} X. Hou, S. Li, and Q. Zheng, Expressing infinite matrices over rings as products of involutions, \emph{Linear Algebra Appl.} \textbf{532} (2017) 257--265.
\bibitem{ofarrell2015} A. G. O'Farrell and I. Short, \emph{Reversibility in Dynamics and Group Theory}, London Mathematical Society Lecture Note Series, Vol. 416 (Cambridge University Press, Cambridge, 2015).
\bibitem{pazzis2021} C. de Seguins Pazzis, Products of involutions of an infinite-dimensional vector space, \emph{Canad. J. Math.} \textbf{73} (2021) 195--220.
\bibitem{pazzis2022} C. de Seguins Pazzis, Sums and products of pairs of quadratic endomorphisms of a countable-dimensional vector space, \emph{Linear Multilinear Algebra} \textbf{70} (2022) 6718--6752; arXiv:2107.04273.
\bibitem{pazzis2024} C. de Seguins Pazzis, A note on bireflectional elements of an algebra, \emph{Comm. Algebra} \textbf{52} (2024) 528--531.
\bibitem{slowik2013} R. S\l{}owik, Expressing infinite matrices as products of involutions, \emph{Linear Algebra Appl.} \textbf{438} (2013) 399--404.
\bibitem{slowik2026} R. S\l{}owik and T. Lohan, On strongly reversible elements in the Riordan group, \emph{Linear Multilinear Algebra} (2026), online publication, doi:10.1080/03081087.2026.2748869.
\bibitem{wonenburger1966} M. J. Wonenburger, Transformations which are products of two involutions, \emph{J. Math. Mech.} \textbf{16} (1966) 327--338.
\end{thebibliography}
\end{document}
